\thispagestyle{empty}
\begin{center}
\vspace*{5.5cm}
{\fontsize{30}{36}\selectfont\bfseries Principia SCADA\par}
\vspace{6pt}
{\bfseries The Principles of the SCADA Technician\par}
\vspace{5.5cm}
{\large Grant Walker\par}
\vspace{4pt}
19 November 2026\par
\end{center}
\clearpage
\renewcommand{\curchap}{Contents}
{\LARGE\bfseries Contents\par}\vspace{4.8pt}\hrule height 0.4pt\vspace{12pt}
{\setlength{\parskip}{0pt}\tocentry{Preface}{3}
\tocpart{Part I. Foundations}
\tocentry{Chapter 1. What SCADA Is}{5}
\tocentry{Chapter 2. Signals and Instruments}{8}
\tocentry{Chapter 3. Points, Quality and Time}{11}
\tocpart{Part II. Communications}
\tocentry{Chapter 4. Field Protocols and Telemetry Links}{14}
\tocentry{Chapter 5. MQTT and Sparkplug B}{18}
\tocentry{Chapter 6. OPC, Gateways and Edge Systems}{22}
\tocentry{Chapter 7. IEC 61850 and Substations}{25}
\tocentry{Chapter 8. Industrial Networks, Timing and Redundancy}{27}
\tocpart{Part III. The SCADA System}
\tocentry{Chapter 9. Servers, Scale and Redundancy}{30}
\tocentry{Chapter 10. HMI Design and Alarm Management}{33}
\tocentry{Chapter 11. Historians and the Unified Namespace}{38}
\tocentry{Chapter 12. Process Control: PID to APC}{42}
\tocentry{Chapter 13. Scripting, APIs and Configuration as Code}{45}
\tocentry{Chapter 14. Virtual PLCs and AI}{49}
\tocpart{Part IV. Security}
\tocentry{Chapter 15. Cryptography for Technicians}{52}
\tocentry{Chapter 16. Securing the Protocols}{55}
\tocentry{Chapter 17. IEC 62443 and Secure Architecture}{58}
\tocentry{Chapter 18. Real Attacks and the Kill Chain}{61}
\tocentry{Chapter 19. Patching, Vulnerabilities and the Law}{64}
\tocpart{Part V. The Working Technician}
\tocentry{Chapter 20. Commissioning, Testing and Fault-Finding}{67}
\tocpart{Appendices}
\tocentry{Appendix A. Key Formulas}{70}
\tocentry{Appendix B. Study Plan and Self-Test}{73}
\tocentry{Appendix C. Glossary and Sources}{76}
\tocentry{Appendix D. Questions to Ask a Vendor}{79}
}
\chap{Preface}
This is a textbook for becoming an excellent SCADA technician. It covers what a technician needs to understand to install, commission, maintain and fault-find a supervisory control system, whichever vendor's products are on site: signals and instruments, protocols and networks, servers and displays, historians and control loops, security, and the method of testing and fault-finding that ties them together.

The book grew from three sets of notes: notes from a Schneider Electric Geo SCADA conference in November 2026, a briefing on eleven advanced topics prepared for the flight there, and the design notes of SecureLink, a hardened MQTT broker for plant networks. Their topics have been taken apart and rebuilt as general principles. Products such as Geo SCADA, AVEVA PI, AVEVA OMI, Kepware and SecureLink appear only as worked examples, marked \emph{In practice}. Learn the idea once; each new product is then a vocabulary exercise.

The method follows the old natural philosophers, which is why the book is called a Principia. A \emph{Definition} fixes the meaning of a word. A \emph{Law} states a relationship that always holds. A \emph{Rule} is good practice learned the hard way. A \emph{Proposition} is a result that follows from the laws. A \emph{Worked Example} applies them to a job you could meet on site, and \emph{On site} notes say how it is done in the field. Each chapter ends with exercises; answers are in brackets.

\sect{What you need before you start}
\begin{itemize}
\item Ohm's law and basic circuit theory. You have this already as an electrician.
\item Arithmetic, percentages and powers. Where the book uses a formula, it is one a technician actually uses: scaling a signal, checking a loop, sizing a link, judging availability.
\item A laptop with Python 3, an MQTT client, a generic OPC UA client and Wireshark, for the practical tasks in Appendix B.
\end{itemize}

\sect{How the book is arranged}
\begin{tbl}{P{0.28\textwidth}L}
Part & Chapters\\\midrule
I. Foundations & 1 What SCADA is; 2 Signals and instruments; 3 Points, quality and time\\
II. Communications & 4 Field protocols; 5 MQTT and Sparkplug B; 6 OPC, gateways and edge; 7 IEC 61850; 8 Networks, timing and redundancy\\
III. The SCADA system & 9 Servers; 10 HMI and alarms; 11 Historians and the UNS; 12 Process control; 13 Scripting and configuration as code; 14 Virtual PLCs and AI\\
IV. Security & 15 Cryptography; 16 Securing the protocols; 17 Accounts, access and audit; 18 IEC 62443 and architecture; 19 Real attacks; 20 Patching and the law\\
V. The working technician & 21 Commissioning, testing and fault-finding\\
\end{tbl}
Short on time? Read Chapters 1, 3, 10, 18 and 21, then Appendix D.

\sect{A note on accuracy}
Protocol rules and mathematics come from the published standards named in each chapter. Product details, dates and legal facts were checked against the sources in Appendix C on 16 and 19 November 2026; some product features had little public documentation yet, and are marked \emph{verify}. Standards are revised and laws change: confirm specifics against the current primary source before relying on them in a formal setting.
